\section{Training Arms That Differ in One Respect}
\label{sec:arms}

All arms fine-tune Qwen3-4B \citep{yang2025qwen3} with thinking disabled. The hyperparameters are identical across arms (learning rate $10^{-5}$, one epoch, global batch 64, maximum length 12,288 tokens). The loss covers the agent's own tokens, that is the note lines and the actions. Observations are never supervised.

Each arm changes one thing relative to the reference arm \arm{N100}.\footnote{Arm names follow the registered protocol. \arm{N100} has fully reliable notes, \arm{NQ}$q$ turns $q$\% of the real notes into foils, the suffix \arm{op} marks a supervised opener, \arm{M} marks a masked action loss, and \arm{SA50} makes 50\% of search actions follow the standard prior (\arm{v2} is the protocol's version label).} Table~\ref{tab:arms} summarises them. We write $S_{\text{data}}$ for the share of real-form notes in the training data whose top type is the receptacle type the teacher visits next.

\begin{table}[t]
\centering\small
\begin{tabular}{@{}lp{0.58\columnwidth}c@{}}
\toprule
Arm & Change relative to \arm{N100} & $S_{\text{data}}$ \\
\midrule
\arm{N100} & none; every note is followed & 1.00 \\
\arm{NQ50op} & half of the real notes are foils; actions unchanged & 0.51 \\
\arm{NQ25op} & a quarter of the real notes are foils & 0.75 \\
\arm{NQ50} & as \arm{NQ50op}, whole foil note unsupervised & 0.51 \\
\arm{SA50v2} & half of all search actions follow the standard prior & 0.66 \\
\arm{N100M} & no loss on actions taken without a note & 1.00 \\
\bottomrule
\end{tabular}
\caption{Training arms. The primary arms \arm{N100} and \arm{NQ50op} have eight seeds each; the other arms have three. A seed fixes both the random draws used to build the data and the data order.}
\label{tab:arms}
\end{table}

\paragraph{\arm{N100}.}
Trajectories come from T$^+$. Every real note states the convention, and the next action visits its top type. Note and convention never disagree, so the data cannot tell a learner whether the action follows the note or follows what the weights know about the object.

\paragraph{Contradicted notes: \arm{NQ50op}, \arm{NQ25op}, \arm{NQ50}.}
These arms keep \arm{N100}'s trajectories and actions token for token. At a fraction $q$ of real-note decision points ($q=0.5$ or $0.25$), the note is replaced by a \emph{foil}, which is the convention ranking of another learnable class of the same scene type, written in the same template. The action still visits the true convention.

In \arm{NQ50op} and \arm{NQ25op} the foil's receptacle names are excluded from the loss, so the model never learns to write a foil. The opener and template of the foil remain supervised. As a result, the share of decision points at which writing a note is supervised is exactly that of \arm{N100}: 0.69--0.71 depending on the draw. \arm{NQ50} excludes the whole foil note from the loss, which lowers that share to 0.39. Comparing \arm{NQ50} with \arm{NQ50op} checks whether writing fewer notes changes the result.

\paragraph{Action-side unreliability: \arm{SA50v2}.}
Notes are always correct, but at every fallback decision point a coin decides whether the teacher takes the convention-following step or the standard-prior step. This happens whether or not a note is written, so the presence of a note does not predict a deviation. In the data the deviation rate is 0.50 with and without a note, and a $\chi^2$ test of independence does not reject at any position or draw.

Standard-prior steps lengthen the search. The arm is therefore subsampled to the same number of fallback decision points as \arm{N100}, about 6,600 of 8,148 trajectories.

\paragraph{No action loss without a note: \arm{N100M}.}
\arm{N100M} is \arm{N100} with the action loss removed at decision points that carry no note (30\%). The model is never trained to act without a note. Section~\ref{sec:related} compares this mask with the action-loss masking of \citet{liu2026cotgains}.
